\documentclass[12pt]{article}
\usepackage[utf8]{inputenc}
\usepackage[T1]{fontenc}
\usepackage{hyperref}
\title{Pom Language Resources for K-12 Curriculum Development}
\author{Compiled by Claude Sonnet 5.5 (claude-sonnet-5-5) for Andr\'as Kornai}
\date{}
\begin{document}
\maketitle

\section{Introduction}
Pom (ISO 639-3 pmo; Glottocode pomm1237) is an Austronesian language of the Yapen area in Cenderawasih Bay, Papua, Indonesia, with about 2,000 speakers. It is poorly documented; the main resource is an online talking dictionary.

\section{Primary Educational Resources}
No primers, textbooks or school programmes were found.

\section{Linguistic and Reference Materials}
\begin{itemize}
\item Pom Dictionary (talking dictionary, recorded 2016) by Claire Benham-Chandler and Emily Gasser, about page: \url{https://talkingdictionary.swarthmore.edu/pom/about}
\item Wikipedia, Pom language: \url{https://en.wikipedia.org/wiki/Pom_language}
\item CrossGram entry for Pom: \url{https://crossgram.clld.org/languages/pomm1237}
\item Numeral data for Pom, MPI-EVA: \url{https://lingweb.eva.mpg.de/channumerals/Pom.htm}
\end{itemize}

\section{Religious and Community Resources}
\begin{itemize}
\item Joshua Project, Pom in Indonesia (people-group profile; its classification differs from Wikipedia's): \url{https://joshuaproject.net/people_groups/14435}
\end{itemize}

\section{Digital and Audio Resources}
The talking dictionary above is the only digital and audio resource found.

\section{Curriculum Development Context}
The talking dictionary offers a starting lexicon with audio for vocabulary materials. A fuller curriculum would need grammar description, orthography agreement and teacher materials, none of which turned up in search.
\end{document}